% Ad Atticum 11.20 (ad-atticum-11-20)
% Source: https://raw.githubusercontent.com/PerseusDL/canonical-latinLit/master/data/phi0474/phi057/phi0474.phi057.perseus-lat1.xml
% Fetched by scripts/fetch_latin.py

% Dateline (Perseus): Scr. Brundisi xvi K. Sept. a. 707 (47).

\ciceroLetterOpener{CICERO ATTICO salutem}

\ciceroSection{1} xvii K. Septembris venerat die xxviii Seleucea Pieria C. Treboni libertus qui se Antiocheae diceret apud Caesarem vidisse Quintum filium cum Hirtio; eos de Quinto quae voluissent impetrasse nullo quidem negotio. quod ego magis gauderem si ista nobis impetrata quicquam ad spem explorati haberent. sed et alia timenda sunt ab aliis Quintis que, et ab hoc ipso quae dantur, ut a domino, rursus in eiusdem sunt potestate. etiam Sallustio ignovit.

\ciceroSection{2} omnino dicitur nemini negare; quod ipsum est suspectum, notionem eius differri. M. Gallius Q. f. mancipia Sallustio reddidit. is venit ut legiones in Siciliam traduceret. eo protinus iturum Caesarem Patris. quod si faciet ego, quod ante mallem, aliquo propius accedam. tuas litteras ad eas quibus a te proxime consilium petivi vehementer exspecto. vale . xvi Kal. Septembris.
